\chapter{Results}\label{chap:results}

This chapter first reports the cost of the runs (Section~\ref{sec:res-setup}) and retrieval quality
(Section~\ref{sec:res-retrieval}). It then reports answer quality and hallucination for all conditions
(Section~\ref{sec:res-main}), the reliability of the judge (Section~\ref{sec:res-judge}), the abstention
rule (Section~\ref{sec:res-abstention}) and an analysis of errors (Section~\ref{sec:res-errors}).
Section~\ref{sec:res-discussion} answers the research questions.

\section{Setup}\label{sec:res-setup}
Table~\ref{tab:cost} shows the mean size and duration of one generation call. Passages make a call longer:
the closed-book calls are the shortest for both models. Every call ended normally, so no answer was cut off
at the token limit. Each hallucination rate below is a share of the non-abstained answers of at most
\val{data.eval_questions} questions, so the confidence intervals are wide and small differences should not be
over-read.

\begin{table}[htbp]
\centering\small
\caption{Generation cost per question: mean total tokens and mean latency of one call, and the share
of calls that ended normally (\texttt{finish\_reason} is \texttt{stop}).}\label{tab:cost}
\begin{tabular}{@{}llrrrr@{}}
\toprule
Model & Condition & Answers & Tokens & Latency (s) & Stopped normally \\
\midrule
gpt-oss-20b & Closed-book & 100 & 273 & 1.93 & 1.000 \\
gpt-oss-20b & BM25 & 100 & 757 & 4.61 & 1.000 \\
gpt-oss-20b & Dense & 100 & 803 & 6.01 & 1.000 \\
gpt-oss-20b & Hybrid & 100 & 814 & 5.94 & 1.000 \\
gpt-oss-20b & Hybrid + rerank & 100 & 785 & 8.70 & 1.000 \\
gpt-oss-20b & Oracle & 100 & 375 & 7.16 & 1.000 \\
\midrule
qwen3.8-27b & Closed-book & 100 & 65 & 1.57 & 1.000 \\
qwen3.8-27b & BM25 & 100 & 693 & 5.33 & 1.000 \\
qwen3.8-27b & Dense & 100 & 743 & 6.33 & 1.000 \\
qwen3.8-27b & Hybrid & 100 & 735 & 5.85 & 1.000 \\
qwen3.8-27b & Hybrid + rerank & 100 & 723 & 5.62 & 1.000 \\
qwen3.8-27b & Oracle & 100 & 305 & 2.71 & 1.000 \\
\bottomrule
\end{tabular}

\end{table}

\section{Retrieval Quality}\label{sec:res-retrieval}

\begin{table}[htbp]
\centering\small
\caption{Retrieval quality on all \val{data.questions} questions and on the first
\val{data.eval_questions}, which are the ones used for generation and judging.}\label{tab:retrieval}
\begin{tabular}{@{}lrrrrrr@{}}
\toprule
 & \multicolumn{3}{c}{All questions} & \multicolumn{3}{c}{Questions used for generation} \\
Retriever & Recall@5 & Both@5 & MRR & Recall@5 & Both@5 & MRR \\
\midrule
BM25 & 0.696 & 0.428 & 0.865 & 0.720 & 0.470 & 0.874 \\
Dense & 0.826 & 0.666 & 0.918 & 0.820 & 0.640 & 0.919 \\
Hybrid & 0.812 & 0.630 & 0.916 & 0.835 & 0.670 & 0.900 \\
Hybrid + rerank & 0.834 & 0.670 & 0.948 & 0.855 & 0.710 & 0.939 \\
\bottomrule
\end{tabular}

\end{table}

\begin{figure}[htbp]
\centering
\resultfigure{fig_retrieval.png}{0.72\textwidth}
\caption{Share of the \val{data.questions} questions with both gold passages in the top $k$
passages.}\label{fig:retrieval}
\end{figure}

Table~\ref{tab:retrieval} and Fig.~\ref{fig:retrieval} show the results. Finding one relevant passage is much
easier than finding both. MRR is high for every retriever (\val{ret.bm25.mrr.all} to
\val{ret.hybrid_rerank.mrr.all} on all questions), yet hybrid + rerank, the best retriever, finds both gold
passages in the top five for only \val{ret.hybrid_rerank.both5.all} of the questions, and BM25 for
\val{ret.bm25.both5.all}.

Two points matter for what follows. First, BM25 is clearly the weakest, so the conditions span a real
range of retrieval quality, which RQ1 needs. Second, the \val{data.eval_questions} questions used for
generation are not a perfect sample of the full set, and the order of dense and hybrid differs between the two:
on all questions dense has the higher Both@5 (\val{ret.dense.both5.all} against \val{ret.hybrid.both5.all}), while
on the first \val{data.eval_questions} hybrid has it (\val{ret.hybrid.both5.eval} against \val{ret.dense.both5.eval}).
No ranking of dense against hybrid is drawn from a difference this small.

\section{Answer Quality and Hallucination}\label{sec:res-main}

\begin{table}[htbp]
\centering\small
\caption{Main results with judge \texttt{gemma-4-31b-it}. Abst.: answers the model itself declined. $n$: judged
answers. Hall.: hallucination rate with its 95\,\% bootstrap interval. Contr., Unsup.: shares of contradicted and
unsupported answers among the judged answers. Both@5 refers to the same \val{data.eval_questions} questions.}\label{tab:main}
\resizebox{\textwidth}{!}{\begin{tabular}{@{}llrrrrrrrr@{}}
\toprule
Model & Condition & EM & F1 & Abst.\,\% & $n$ & Hall.\,\% [95\,\% CI] & Contr.\,\% & Unsup.\,\% & Both@5 \\
\midrule
gpt-oss-20b & Closed-book & 0.19 & 0.29 & 0 & -- & -- & -- & -- & -- \\
gpt-oss-20b & BM25 & 0.42 & 0.57 & 11 & 88 & 34.1 [25.0, 43.2] & 1.1 & 33.0 & 0.47 \\
gpt-oss-20b & Dense & 0.52 & 0.66 & 8 & 90 & 21.1 [13.3, 30.0] & 1.1 & 20.0 & 0.64 \\
gpt-oss-20b & Hybrid & 0.51 & 0.67 & 4 & 93 & 23.7 [15.1, 32.3] & 1.1 & 22.6 & 0.67 \\
gpt-oss-20b & Hybrid + rerank & 0.51 & 0.67 & 2 & 96 & 25.0 [16.7, 34.4] & 2.1 & 22.9 & 0.71 \\
gpt-oss-20b & Oracle & 0.65 & 0.79 & 1 & 98 & 8.2 [3.1, 13.3] & 3.1 & 5.1 & 1.00 \\
\midrule
qwen3.8-27b & Closed-book & 0.21 & 0.28 & 0 & -- & -- & -- & -- & -- \\
qwen3.8-27b & BM25 & 0.54 & 0.66 & 1 & 97 & 42.3 [32.0, 51.6] & 3.1 & 39.2 & 0.47 \\
qwen3.8-27b & Dense & 0.66 & 0.76 & 1 & 96 & 28.1 [19.8, 37.5] & 1.0 & 27.1 & 0.64 \\
qwen3.8-27b & Hybrid & 0.61 & 0.73 & 1 & 97 & 29.9 [21.6, 39.2] & 4.1 & 25.8 & 0.67 \\
qwen3.8-27b & Hybrid + rerank & 0.63 & 0.74 & 2 & 96 & 25.0 [16.7, 34.4] & 2.1 & 22.9 & 0.71 \\
qwen3.8-27b & Oracle & 0.73 & 0.84 & 1 & 98 & 6.1 [2.0, 11.2] & 1.0 & 5.1 & 1.00 \\
\bottomrule
\end{tabular}
}
\end{table}

\begin{figure}[htbp]
\centering
\resultfigure{fig_main.png}{\textwidth}
\caption{Hallucination rate by condition and model with 95\,\% bootstrap intervals. The Both@5 values under the
conditions refer to the same \val{data.eval_questions} questions.}\label{fig:main}
\end{figure}

Table~\ref{tab:main} and Fig.~\ref{fig:main} give the main results. Retrieval helps accuracy: EM rises from 0.19 and
0.21 closed-book (\texttt{gpt-oss-20b} and \texttt{qwen3.8-27b}) to at best 0.52 and 0.66 with a retriever, and to
0.65 and 0.73 with the oracle. Hallucination is highest for BM25 (34.1\,\% and 42.3\,\%) and lowest for the oracle
(8.2\,\% and 6.1\,\%). The intervals of the four retrievers overlap for both models, including BM25 against dense,
so the retrievers are not ranked. Nor does a better Both@5 translate neatly at this scale: hybrid + rerank has the
highest Both@5 on these questions, but for \texttt{gpt-oss-20b} its hallucination rate is not below that of dense.

Almost all hallucinations are unsupported, not contradicted. Of the 949 judged answers, 230 are hallucinated: 211
unsupported and 19 contradicted. The models rarely state something that the passages conflict with; they more often
state something that the passages do not establish. Of the 968 non-abstained answers with passages, 19 (2.0\,\%) got
no readable verdict from the main judge and are excluded.

\subsection{Effect of the Gold Passages Found}

\begin{table}[htbp]
\centering\small
\caption{Hallucination by the number of gold passages among the five passages shown to the model (the four
retrievers pooled; the oracle separately). No rate is given for 0 of 2, which rests on one and two answers. CI
(answers) resamples answers and is too narrow for pooled rows; CI (questions) resamples whole
questions.}\label{tab:gold}
\begin{tabular}{@{}llrrrrr@{}}
\toprule
Model & Gold found & $n$ & Hall. & Rate\,\% & CI (answers) & CI (questions) \\
\midrule
gpt-oss-20b & 0 of 2 & 1 & 1 & -- & -- & -- \\
gpt-oss-20b & 1 of 2 & 122 & 77 & 63.1 & [54.9, 71.3] & [48.4, 77.0] \\
gpt-oss-20b & 2 of 2 & 244 & 17 & 7.0 & [3.7, 10.2] & [2.7, 12.0] \\
gpt-oss-20b & Oracle & 98 & 8 & 8.2 & [3.1, 13.3] & -- \\
\midrule
qwen3.8-27b & 0 of 2 & 2 & 2 & -- & -- & -- \\
qwen3.8-27b & 1 of 2 & 139 & 108 & 77.7 & [70.5, 84.2] & [65.7, 88.3] \\
qwen3.8-27b & 2 of 2 & 245 & 11 & 4.5 & [2.0, 7.3] & [0.8, 9.3] \\
qwen3.8-27b & Oracle & 98 & 6 & 6.1 & [2.0, 11.2] & -- \\
\bottomrule
\end{tabular}

\end{table}

\begin{figure}[htbp]
\centering
\resultfigure{fig_gold_found.png}{0.62\textwidth}
\caption{Hallucination rate against the number of gold passages in the top five, with bootstrap intervals that
resample answers. The bars for 0 of 2 rest on one and two answers and should not be interpreted.}\label{fig:gold}
\end{figure}

Table~\ref{tab:gold} and Fig.~\ref{fig:gold} hold the key result for RQ1. With both gold passages in the top five,
7.0\,\% and 4.5\,\% of the answers are hallucinated, in line with the oracle. With only one, the rates are 63.1\,\%
and 77.7\,\%. Retrieval found neither passage in only three cases, too few to interpret.

Because the table pools four retrievers, one question can appear up to four times. Resampling whole questions widens
the intervals but does not close the gap (Table~\ref{tab:gold}). Two further checks address whether questions with a
missed hop are simply harder. First, some of them are: under the oracle, questions that had a one-gold retrieval reach
15.7\,\% and 8.5\,\% hallucination, against 5.1\,\% and 3.8\,\% for questions that had a two-gold retrieval. That
difference is far smaller than the one in Table~\ref{tab:gold}. Second, on questions where the retrievers disagree,
each question can be compared with itself. On these 32 and 40 questions, hallucination is 64.1\,\% and 78.3\,\% with
one gold passage and 10.4\,\% and 5.0\,\% with two. This remains an observation and not a controlled experiment, but
difficulty alone does not explain the effect.

The result is partly built into the definition: with a hop missing, the passages can rarely support the answer. What
the data add is how the models behave in that state. The prompt offered no explicit way to decline, but it did tell
the models to use only the context, and they rarely refused: self-abstention never exceeds 11\,\% in any condition and
is 1--2\,\% for \texttt{qwen3.8-27b} when it is given passages. A missing hop usually leads to an answer anyway.

\subsection{Correctness against Faithfulness}

\begin{table}[htbp]
\centering\small
\caption{Counts of correct (F1 $\geq$ 0.5) or incorrect answers that were labelled supported or
hallucinated.}\label{tab:correctness}
\begin{tabular}{@{}llrrrr@{}}
\toprule
 &  & \multicolumn{2}{c}{Correct} & \multicolumn{2}{c}{Incorrect} \\
Model & Condition & Supp. & Hall. & Supp. & Hall. \\
\midrule
gpt-oss-20b & BM25 & 44 & 17 & 14 & 13 \\
gpt-oss-20b & Dense & 58 & 10 & 13 & 9 \\
gpt-oss-20b & Hybrid & 58 & 12 & 13 & 10 \\
gpt-oss-20b & Hybrid + rerank & 56 & 14 & 16 & 10 \\
gpt-oss-20b & Oracle & 76 & 4 & 14 & 4 \\
\midrule
qwen3.8-27b & BM25 & 49 & 21 & 7 & 20 \\
qwen3.8-27b & Dense & 63 & 14 & 6 & 13 \\
qwen3.8-27b & Hybrid & 61 & 14 & 7 & 15 \\
qwen3.8-27b & Hybrid + rerank & 63 & 13 & 9 & 11 \\
qwen3.8-27b & Oracle & 81 & 5 & 11 & 1 \\
\bottomrule
\end{tabular}

\end{table}

Faithfulness and correctness come apart (Table~\ref{tab:correctness}). Of the 230 hallucinated answers, 124 (54\,\%)
are nevertheless correct by the F1 threshold (57 of 103 for \texttt{gpt-oss-20b}, 67 of 127 for \texttt{qwen3.8-27b}).
This fits the models filling a missing hop from what they already know, or guessing well, although these counts
cannot separate the two. In the other direction, 110 of 719 supported answers (15\,\%) are incorrect. A check for
faithfulness cannot replace a check for accuracy, nor the reverse.

\subsection{Question Type}

\begin{table}[htbp]
\centering\small
\caption{EM and hallucination rate for bridge and comparison questions. Only
\val{data.eval_comparison} of the \val{data.eval_questions} questions are comparison
questions.}\label{tab:qtype}
\begin{tabular}{@{}llrrrr@{}}
\toprule
 &  & \multicolumn{2}{c}{EM} & \multicolumn{2}{c}{Hallucination rate (\%)} \\
Model & Condition & Bridge & Comp. & Bridge & Comp. \\
\midrule
gpt-oss-20b & BM25 & 0.37 & 0.62 & 32.4 & 40.0 \\
gpt-oss-20b & Dense & 0.47 & 0.71 & 26.1 & 4.8 \\
gpt-oss-20b & Hybrid & 0.46 & 0.71 & 27.8 & 9.5 \\
gpt-oss-20b & Hybrid + rerank & 0.47 & 0.67 & 30.7 & 4.8 \\
gpt-oss-20b & Oracle & 0.62 & 0.76 & 10.3 & 0.0 \\
\midrule
qwen3.8-27b & BM25 & 0.49 & 0.71 & 43.4 & 38.1 \\
qwen3.8-27b & Dense & 0.65 & 0.71 & 36.0 & 0.0 \\
qwen3.8-27b & Hybrid & 0.58 & 0.71 & 36.8 & 4.8 \\
qwen3.8-27b & Hybrid + rerank & 0.62 & 0.67 & 30.3 & 5.0 \\
qwen3.8-27b & Oracle & 0.73 & 0.71 & 7.7 & 0.0 \\
\bottomrule
\end{tabular}

\end{table}

Pooled across the five conditions with passages, comparison questions hallucinate less than bridge questions
(11.7\,\% against 25.1\,\% for \texttt{gpt-oss-20b}, 9.7\,\% against 30.7\,\% for \texttt{qwen3.8-27b}). This does not
hold for BM25, where comparison questions reach 40.0\,\% and 38.1\,\% (Table~\ref{tab:qtype}). There are only
\val{data.eval_comparison} comparison questions per condition, so this is suggestive only.

\section{Reliability of the Judge}\label{sec:res-judge}

\begin{table}[htbp]
\centering\small
\caption{Agreement between \texttt{gemma-4-31b-it} and \texttt{openai/gpt-oss-120b}.}\label{tab:agreement}
\begin{tabular}{@{}lrrrrr@{}}
\toprule
 & $n$ & Agree 3 labels\,\% & $\kappa$ 3 labels & Agree binary\,\% & $\kappa$ binary \\
\midrule
The two judges & 663 & 90.5 & 0.77 & 92.3 & 0.81 \\
\bottomrule
\end{tabular}

\end{table}

\begin{table}[htbp]
\centering\small
\caption{Labels of the two judges (rows: main judge; columns: second judge).}\label{tab:confusion}
\begin{tabular}{@{}lrrr@{}}
\toprule
 & supported & contradicted & unsupported \\
\midrule
supported & 454 & 4 & 13 \\
contradicted & 2 & 5 & 7 \\
unsupported & 32 & 5 & 141 \\
\bottomrule
\end{tabular}

\end{table}

\begin{table}[htbp]
\centering\small
\caption{Agreement of each judge with the manual labels of the author. The sample is balanced by the main judge's
verdict, so these values are not an accuracy estimate for all answers.}\label{tab:human}
\begin{tabular}{@{}lrrrrr@{}}
\toprule
Compared with & $n$ & Agree 3 labels\,\% & $\kappa$ 3 labels & Agree binary\,\% & $\kappa$ binary \\
\midrule
gemma-4-31b-it & 100 & 73.0 & 0.48 & 77.0 & 0.54 \\
gpt-oss-120b & 79 & 77.2 & 0.52 & 78.5 & 0.54 \\
\bottomrule
\end{tabular}

\end{table}

\textbf{Between the judges.} The two judges agree on 92.3\,\% of 663 shared answers with binary labels ($\kappa=0.81$)
and on 90.5\,\% with three labels ($\kappa=0.77$; Tables~\ref{tab:agreement} and~\ref{tab:confusion}). Agreement between
models does not show that either is right, and the second judge shares a family with one generator.

\textbf{Human check.} The author labelled all 100 answers of the shuffled sheet (Table~\ref{tab:human}). Agreement
with the main judge is 77.0\,\% with binary labels ($\kappa=0.54$) and 73.0\,\% with three labels ($\kappa=0.48$), which
is moderate and clearly lower than between the two judges, although the balanced sample means the values are not
directly comparable. The disagreements are one-sided: of 23 binary disagreements with the main judge, the judge called
20 answers hallucinated where the author called them supported, and only 3 the other way round. Reweighting the
sample to the observed rate gives about 19.1\,\% hallucination overall against the judge's 24.2\,\%. The share
of hallucinated labels is similar in the first 80 answers (26) and the last 20 (7).

There are two readings of this gap, and the data cannot separate them. The judge may be too strict: it sometimes
faults an answer for wording in the question. The author may also be too lenient: in 10 of those 20 cases, both judges
called the answer hallucinated. The absolute rates in Table~\ref{tab:main} are therefore
uncertain and possibly too high. The main result does not depend on this: using the author's labels, 26 of 42 answers
with one gold passage are hallucinated, against 3 of 42 with both (about 54\,\% and 4\,\% after reweighting within each
group). These subsets are small and balanced, so they support the direction and size of the gap, not exact values.

\section{Abstention}\label{sec:res-abstention}

\begin{table}[htbp]
\centering\footnotesize
\caption{Abstention on held-out answers, as the mean over 200 random development/test splits (about 49 answers per test
half). AUROC is computed on all answers of the model. Cover.: coverage. Hall., Acc.: hallucination rate and accuracy,
before and after abstention; SD over the overlapping splits in brackets, which is not a confidence interval. Rem.:
hallucinations withheld per test half. Lost: correct answers withheld per test half. The ``before'' values are averaged
over test halves, so they differ slightly from Table~\ref{tab:main}.}\label{tab:abstention}
\resizebox{\textwidth}{!}{\begin{tabular}{@{}lllrrrrrrrr@{}}
\toprule
 &  &  &  &  & \multicolumn{2}{c}{Hall.\,\%} &  &  & \multicolumn{2}{c}{Acc.\,\%} \\
Retriever & Model & Signal & AUROC & Cover.\,\% & Before & After (SD) & Rem. & Lost & Before & After \\
\midrule
Hybrid + rerank & gpt-oss-20b & top1 & 0.50 & 95.9 & 25.4 & 25.9 (4.1) & 0.0 & 0.4 & 71.4 & 70.6 \\
Hybrid + rerank & gpt-oss-20b & top2 & 0.57 & 93.9 & 25.4 & 24.3 (4.5) & 1.0 & 1.5 & 71.4 & 68.3 \\
Hybrid + rerank & gpt-oss-20b & answer\_overlap & 0.75 & 81.8 & 25.4 & 17.1 (4.9) & 5.3 & 3.2 & 71.4 & 65.0 \\
Hybrid + rerank & qwen3.8-27b & top1 & 0.48 & 96.3 & 25.2 & 25.2 (4.5) & 0.2 & 0.2 & 77.3 & 76.9 \\
Hybrid + rerank & qwen3.8-27b & top2 & 0.54 & 93.8 & 25.2 & 25.0 (4.4) & 0.6 & 1.4 & 77.3 & 74.4 \\
Hybrid + rerank & qwen3.8-27b & answer\_overlap & 0.60 & 91.5 & 25.2 & 22.3 (5.0) & 2.0 & 1.5 & 77.3 & 74.3 \\
\midrule
Dense & gpt-oss-20b & top1 & 0.54 & 89.1 & 21.3 & 21.4 (4.5) & 0.2 & 1.1 & 69.6 & 67.5 \\
Dense & gpt-oss-20b & top2 & 0.68 & 86.0 & 21.3 & 19.2 (5.3) & 1.4 & 1.7 & 69.6 & 66.2 \\
Dense & gpt-oss-20b & answer\_overlap & 0.77 & 74.9 & 21.3 & 11.0 (4.9) & 5.5 & 2.8 & 69.6 & 63.9 \\
Dense & qwen3.8-27b & top1 & 0.46 & 98.0 & 28.2 & 28.4 (4.7) & 0.0 & 0.5 & 79.0 & 78.1 \\
Dense & qwen3.8-27b & top2 & 0.53 & 96.2 & 28.2 & 27.7 (4.9) & 0.6 & 1.1 & 79.0 & 76.7 \\
Dense & qwen3.8-27b & answer\_overlap & 0.63 & 90.2 & 28.2 & 23.2 (5.5) & 3.4 & 2.3 & 79.0 & 74.3 \\
\bottomrule
\end{tabular}
}
\end{table}

\begin{figure}[htbp]
\centering
\resultfigure{fig_abstention.png}{\textwidth}
\caption{AUROC of each signal. The vertical axis starts at 0.3; the dashed line marks chance.}\label{fig:auroc}
\end{figure}

\textbf{Retrieval confidence does not work.} The top1 score is near chance everywhere (AUROC 0.46--0.54) and removes
almost no hallucinations (Table~\ref{tab:abstention}, Fig.~\ref{fig:auroc}). The top2 score ranges from 0.53 to 0.68. It
is better in one case (\texttt{gpt-oss-20b} with dense), but the thresholds chosen for it keep most answers and change
hallucination by little. A likely reason, which was not tested directly, is that the score reflects how well a passage
matches the question, and a strong match to one hop says little about whether the other was found.

\textbf{The grounding check helps, unevenly.} The answer\_overlap signal reaches an AUROC of 0.75--0.77 for
\texttt{gpt-oss-20b} and 0.60--0.63 for \texttt{qwen3.8-27b}. Its best case is \texttt{gpt-oss-20b} with dense, where
hallucination falls from 21.3\,\% to 11.0\,\% at 74.9\,\% coverage. The cost is real: per test half it withholds 5.5
hallucinations and 2.8 correct answers, and accuracy falls from 69.6\,\% to 63.9\,\%. For \texttt{qwen3.8-27b} the gain is
much smaller (28.2\,\% to 23.2\,\% with dense).

\section{Error Analysis}\label{sec:res-errors}

\begin{table}[htbp]
\centering\small
\caption{Causes of \val{cfg.n_errors} hallucinated answers from the BM25 and dense conditions, assigned by the
author.}\label{tab:errors}
\begin{tabular}{@{}lr@{}}
\toprule
Cause & Author \\
\midrule
one hop missing & 18 \\
wrong passage retrieved & 8 \\
answer from memory & 0 \\
passage misread & 3 \\
judge error & 1 \\
\bottomrule
\end{tabular}

\end{table}

\textbf{Where the words come from.} For \texttt{qwen3.8-27b}, all words of the answer appear in the context in 91 of 127
hallucinated answers (72\,\%), against 48 of 103 (47\,\%) for \texttt{gpt-oss-20b}. This suggests that Qwen more often
answers with a span from the passages that they do not back, which a lexical check cannot see. It would also explain
why the grounding check works better for the model whose hallucinations more often contain words from outside the
context.

\textbf{Causes.} Table~\ref{tab:errors} shows the causes of \val{cfg.n_errors} hallucinated answers. Causes overlap. The
author assigned each answer one of the five causes shown in the labelling tool.
The author used wrong passage retrieved for 8 answers, all with one gold passage in the top five. The author never chose answer from memory.

A coarser reading is that the author attributes 26 of the 30
answers to evidence missing from the context (one hop missing or wrong passage retrieved). Of the 30 answers, 27 had one gold passage in the top
five, including all 26.

The judge-error count does not fit the human check. The author chose judge error for 1 of the 30 answers, but on the
label sheet called 20 of the 50 answers that the judge called hallucinated supported. Six answers are on both sheets: the
author labelled 3 of them supported, and gave only 1 of those 3 the cause judge error. The cause step showed the judge's
reason, which may have influenced the author, so the author's count of judge errors may be too low.

Three real examples show the pattern. Asked in what year Philip Pullman's book was put into audio, a model answered
``1995'': the passages give that as the publication year and do not mention the audio version (wrong passage retrieved). Asked
what year the subject of \emph{Fatal Vision} was convicted, a model gave the correct ``1979'', a year that appears nowhere
in the passages (one hop missing). Asked what Iqbal F. Qadir was on during an attack, a model answered ``Vice-Admiral'',
although the passage says that he was part of a flotilla (passage misread).

\section{Discussion}\label{sec:res-discussion}

\textbf{RQ1: How does retrieval quality affect hallucination?}
Whether retrieval found every hop matters much more than which retriever was used. With both gold passages in the top
five, 7.0\,\% and 4.5\,\% of the answers are hallucinated; with one, 63.1\,\% and 77.7\,\%. The gap holds when a question is
compared with itself and when the author's labels replace the judge's. Differences between the retrievers are smaller and
their intervals overlap. For multi-hop questions this suggests tracking Both@$k$ alongside MRR, since MRR was high here even
where a hop was often missing. Improving second-hop recall, for example with iterative retrieval, looks like a promising
target, although it was not tested here.

\textbf{RQ2: Can a training-free rule reduce hallucination without losing many correct answers?}
Yes, but not with the signal that pipelines most readily have, and not without a cost. Retrieval scores should not be
assumed to indicate whether an answer is supported: the top1 score is near chance. The grounding check needs only string
operations after generation, but its benefit depends on the model, and it can catch only answers with words from outside
the context. It costs correct answers, partly because 54\,\% of the hallucinated answers are correct by F1: withholding them
improves faithfulness at the expense of accuracy. The rule is training-free but not label-free, because choosing $\tau$
needs about 50 labelled answers.

\textbf{What this means in practice.}
A multi-hop RAG system should measure whether all the evidence was retrieved, not only whether something relevant was. A
high retrieval score is no reason to trust an answer. A simple grounding check is worth trying where a withheld answer is
cheaper than an unsupported one, but its value should be measured for the model in use. Self-RAG and CRAG were not run on
this data, so they are compared here only in what they require: they need fine-tuning, and the rule studied here does not.
The limitations of these conclusions are discussed in Section~\ref{sec:limitations}.
